\section{推論と意思決定}
\label{D1_lesson17}

\subsection{授業内容}
\subsubsection{科目の中でのこのコマの位置づけ}
\begin{tabular}[t]{cccccc}
    記述統計[4] & 確率[4] & モデル[8] & 推測・推定[2] & 意思決定[1] & 実践[2]
\end{tabular}
\subsubsection{概要}
標本統計量をつかって母数を推定する方法として，幅を持たせた推論である区間推定という方法がある。この区間の中に母数があるとはどういうことか，注意して理解する必要がある。また標本統計量を標準化しておくことで，標準正規分布を使って一般的な推定ができる利便性があることを理解する。これらは推定の問題であったが，心理学における多くの場合には推定量そのものよりも，平均因果効果があったのかなかったのか，という判断を下したいという目的があり，そのために検定と呼ばれるモデル比較の考え方が導入される。正規分布など母集団分布におけるモデルはそのままに，パラメータがどういう値であるかについての仮説を戦わせるその作法を学ぶ。
\subsubsection{コマ主題細目(箇条書き)}
\begin{description}
\item[点推定と区間推定]標本平均は母平均の不偏推定量ではあるが，ある値をそのまま母平均であると考えるのは(標準誤差の考え方があるとは言え)やや断定的である。これに対して一定の幅を使って予測する区間推定という方法がある。推定に確率の言葉が入ってくるので不慣れな表現が出てくるが，注意深く解釈する必要がある。
\begin{flushright}
    $\to$ \citet{川端201408}のP.111--116.
\end{flushright}

\item[標準正規分布の利点]標本ごとに頻度主義を描くのではなく，標本統計量を標準化しておくことで標準正規分布に対応させ，例えば95\%区間に対応する$\pm 1.96$という数字を使うことができれば便利である。ここでは標準化の計算についておさらいするとともに，標準正規分布の表の利用の仕方について理解する。
\begin{flushright}
    $\to$ \citet{Yamada2004}のP.86--89.
\end{flushright}

\item[推定と検定]ここまでは平均値などを推定してきたが，推定量を意思決定に応用する方法がある。特定の数値についての複数のモデルを比較し優劣を決めるモデル比較の考え方がこれで，心理学においては特に，帰無仮説と対立仮説という2つのモデルを用いる帰無仮説検定という考え方である。
\begin{flushright}
    $\to$ \citet{川端201408}のP.81--85.
\end{flushright}

\item[統計的帰無仮説検定]統計的帰無仮説検定(Null Hypothesis Significance Test;NHST)は帰無仮説に基づいてモデル比較をする一連の手続きである。この手続きはさまざまな統計量に対して一般化されているため，統計パッケージの中にも含まれており，統計法の専門家でなくても簡単に利用することができる。ただし安易な使い方は誤用，悪用を招くことになるので，慎重にその手続きとロジックを理解しておく必要があるだろう。
\begin{flushright}
    $\to$ 手続きの一般化としては\citet{Yamada2004}のP.108--109　
\end{flushright}

\item[有意水準]モデル比較という観点からは，どうしても「判定基準」を考える必要がある。NHSTの文脈では，それは希少性の指標である有意水準になる。有意水準の定め方は領域ごとに定める任意であり，心理学では一般に5\%が用いられている。このことは，複数のモデルの優劣を判定する数字であって，モデルの強さ，正しさを表現する数字ではないことに注意する。
\begin{flushright}
    $\to$ \citet{Yamada2004}のP.112--115.
\end{flushright}

\end{description}
\subsubsection{キーワード}
\begin{itemize}
\item 区間推定
\item 標準正規分布
\item 帰無仮説検定(NHST)
\item 有意水準
\end{itemize}
\subsection{授業情報}
\paragraph{コマの展開方法}
講義
\paragraph{標準シラバスにおける位置づけ}
科目番号5；心理学統計法；1.心理学で用いられる統計手法；F 推測統計:その考え方
\subsubsection{予習・復習課題}
\paragraph{予習}
母集団と標本の関係，標本から母数を推定できることについて，前回の話の延長線上で進むので十分に復習してから授業に挑むこと。特に限られた情報から未知なるものを推定するという手続きに当たっては，幾重にも過程が重ねられていくので，どのような仮定を置いて進めてきているかを再確認しておくと良いだろう。
\paragraph{復習}
帰無仮説検定については，触れていない心理統計のテキストはないと言っても過言ではないので，どのようなものでも良いから手に取り，その手続きや用語について復習しておくことが望ましい。テキストによっては，分布(ex. $t$検定，$F$検定)や要因計画(ex.対応のあるt検定，分散分析)ごとに分けて書かれていたり，フローチャートのように説明されているものもあるが，本質は母数の推定とモデル比較であることに注意して読み解くと良い。